\begin{theorem}[Геометрический смысл градиента]
	Если $f$ дифференцируема в точке $\v x_0\in\R^n$ и $\grad f(\v x_0)\neq\v0$, то $\frac{\vd f}{\vd\v l}$ максимальна среди всех $\v l,\ |\v l|=1$, для $\v l$ сонаправленного с $\grad f(\v x_0)$, причём в этом случае $\frac{\vd f}{\vd\v l}=|\grad f(\v x_0)|$, и минимальна среди тех же $\v l$ для $\v l$, противоположно направленных к $\grad f(\v x_0)$, причём тогда $\frac{\vd f}{\vd\v l}=-|\grad f(\v x_0)|$
\end{theorem}
\begin{proof}
	Так как от $\v l$ требуется лишь направление, будем считать $|\v l|=1$.\\
	\[\frac{\vd f}{\vd\v l}=|\grad f(\v x_0),\v l|\]
	По неравенству Коши-Буняковского-Шварца имеем:
	\[\frac{\vd f}{\vd\v l}\le|\grad f(\v x_0)|\cdot|\v l|\text{, причём равенство имеет место т. и т. т. к. }\grad f(\v x_0)=\alpha\v l\text{, то есть:}\]
	\[\v l=\pm\frac{\grad f(\v x_0)}{|\grad f(\v x_0)|}\]
	\[\v l=\frac{\grad f(\v x_0)}{|\grad f(\v x_0)|}\Ra\frac{\vd f}{\vd\v l}=\frac{(\grad f(\v x_0),\grad f(\v x_0))}{|\grad f(\v x_0)|}=|\grad f(\v x_0)|\]
\end{proof}
\newpage
\section{Частные производные и дифференциалы высших порядков}
\begin{anote}
	По каким - то причинам, в нумерации Алексея Леонидовича этот параграф стоит под номером 4, хотя он уже был. В целом ладно, но читателю стоит иметь это ввиду. Вероятно, далее нумерация будет сбита.
\end{anote}
\begin{definition}
	Если частная производная $\frac{\vd f}{\vd x_i}$ определена в некоторой окрестности точки $\v x_0$, то её частная производная (если она $\exists$) по переменной $x_j$ в точке $\v x_0$ называется (смешанной при $i\neq j$) \textit{частной производной второго порядка}:
	\[\frac{\vd^2f}{\vd x_j\vd x_i}(\v x_0):=\frac{\vd}{\vd x_j}\left(\frac{\vd f}{\vd x_i}\right)(\v x_0)\]
	Частные производные высших порядков определяются по индукции:
	\[\frac{\vd^mf}{\vd x_{i_m}\vd x_{i_{m-1}}\cdot\vd x_{i_1}}(\v x_0):=\frac{\vd}{\vd x_{i_m}}\left(\frac{\vd^{m-1}f}{\vd x_{i_{m-1}},\dots,\vd x_{i_1}}\right)(\v x_0)\]
\end{definition}
\begin{note}
	Смешанная производная зависит от порядка дифференцирования.
\end{note}
\begin{example}
	\[
	f(x,y)=
	\begin{cases}
		xy\cdot\frac{x^2-y^2}{x^2+y^2},\ (x,y)\neq(0,0)\\
		0,\ (x,y)=(0,0)
	\end{cases}
	\]
	\[\frac{\vd f}{\vd x}(x,y)=y\cdot\frac{x^2-y^2}{x^2+y^2}+xy\cdot\frac{2x(x^2+y^2)-2x(x^2-y^2)}{(x^2+y^2)^2}=y\cdot\frac{x^2-y^2}{x^2+y^2}+\frac{3x^2y^3}{(x^2+y^2)^2},\ (x,y)\neq(0,0)\]
	\[\frac{\vd f}{\vd y}(x,y)=x\cdot\frac{x^2-y^2}{x^2+y^2}-\frac{4x^3y^2}{(x^2+y^2)^2}\]
	\begin{equation*}
		\begin{split}
			\frac{\vd f}{\vd x}(0,0)=\frac{d}{dx}\bigg(f(x,0)\bigg)&\bigg|_{x=0}=0=\frac{\vd f}{\vd y}(0,0)\\
			\frac{\vd^2f}{\vd x\vd y}(0,0)=\frac{d}{dx}\left(\frac{\vd f}{\vd y}(x,0)\right)&\bigg|_{x=0}=\frac{d}{dx}(x)\bigg|_{x=0}=1\\
			\frac{\vd^2f}{\vd y\vd x}(0,0)=\frac{d}{dy}\bigg(\frac{\vd f}{\vd x}(0,y)\bigg)&\bigg|_{y=0}=\frac{d}{dy}(-y)\bigg|_{y=0}=-1
		\end{split}
	\end{equation*}
\end{example}
\begin{theorem}[Шварц]\label{th1}
	Если $\frac{\vd^2f}{\vd x\vd y}$ и $\frac{\vd^2f}{\vd y\vd x}$ $\exists$ в некоторой окрестности точки $(x_0,y_0)$ и непрерывны в ней, то: 
	\[\frac{\vd^2f}{\vd x\vd y}(x_0,y_0)=\frac{\vd^2f}{\vd y\vd x}(x_0,y_0)\]
\end{theorem}
\begin{proof}
	Рассмотрим $\psi(h)=f(x_0+h,y_0+h)-f(x_0+h,y_0)-f(x_0,y_0+h)+f(x_0,y_0)$
	\[\psi(h)=\tr\phi=\phi(x_0+h)-\phi(x_0)\text{, где }\phi(x)=f(x,y_0+h)\]
	Отметим, что $\phi$ дифференцируема в некоторой окрестности точки $x_0$ при $h\ra0$.\\
	По т. Лагранжа $\exists\ \xi_1$ между $x_0$ и $x_0+h$ такой, что: 
	\[\psi(h)=\phi'(\xi_1)\cdot h=\bigg(\frac{\vd f}{\vd x}(\xi_1,y_0+h)-\frac{\vd f}{\vd x}(\xi_1,y_0)\bigg)\cdot h\]
	Аналогично, по т. Лагранжа $\exists\ \eta_1$ между $y_0$ и $y_0+h$.
	\[\text{Получим: }\bigg(\frac{\vd f}{\vd x}(\xi_1,y_0+h)-\frac{\vd f}{\vd x}(\xi_1,y_0)\bigg)\cdot h=\frac{\vd^2f}{\vd y\cdot\vd x}(\xi_1,\eta_1)\cdot h^2\]
	Проделаем аналогичные вещи, но теперь с $y$:
	\[\psi(h)=\tr\mu=\mu(y_0+h)-\mu(y_0)\text{, где }\mu(y)=f(x_0+h,y)-f(x_0,y)\]
	\newpage
	Воспользуемся т. Лагранжа ещё 2 раза, получим $\xi_2$ между $x_0$ и $x_0+h$, а также $\eta_2$ между $y_0$ и $y_0+h$.
	\[\psi(h)=\mu'(\eta_2)\cdot h=\bigg(\frac{\vd f}{\vd y}(x_0+h,\eta_2)-\frac{\vd f}{\vd y}(x_0,\eta_2)\bigg)\cdot h=\frac{\vd^2f}{\vd x\cdot\vd y}(\xi_2,\eta_2)\cdot h^2\]
	Итого имеем:
	\[\frac{\vd^2f}{\vd y\cdot\vd x}(\xi_1,\eta_1)=\frac{\vd^2f}{\vd x\cdot\vd y}(\xi_2,\eta_2),\ h\ra0\]
	Получили требуемое.
\end{proof}
\begin{definition}
	$f$ называется \textit{дважды дифференцируемой} в точке $\v x_0\in\R^n$, если все её частные производные $\frac{\vd f}{\vd x_i},\ i=1,\dots,n$ $\exists$ в окрестности точки $\v x_0$ и дифференцируемы в ней.
\end{definition}
\begin{theorem}[Юнга]
	Если $f(x,y)$ дважды дифференцируема в точке $(x_0,y_0)$, то $\frac{\vd^2f}{\vd x\cdot\vd y}(x_0,y_0)=\frac{\vd^2f}{\vd y\cdot\vd x}(x_0,y_0)$
\end{theorem}
\begin{proof}
	Аналогично теореме{~\eqref{th1}} $\psi(h)=f(x_0+h,y_0+h)-f(x_0,y_0+h)-f(x_0+h,y_0)+f(x_0,y_0)$
	\begin{multline*}
		\psi(h)=\tr\phi=\phi(x_0+h)-\phi(x_0)=\bigg(\frac{\vd f}{\vd x}(\xi,y_0+h)-\frac{\vd f}{\vd x}(\xi,y_0)\bigg)\cdot h=\\
		=\bigg(\frac{\vd f}{\vd x}(\xi,y_0+h)-\frac{\vd f}{\vd x}(x_0,y_0)\bigg)\cdot h-
		\bigg(\frac{\vd f}{\vd x}(\xi,y_0)-\frac{\vd f}{\vd x}(x_0,y_0)\bigg)\cdot h=\text{(следует из определения дифф-ти)}\\
		=\bigg(\frac{\vd^2f}{\vd x^2}(x_0,y_0)(\xi-x_0)+\frac{\vd^2f}{\vd y\cdot\vd x}(x_0,y_0)\cdot h+\sigma\big|(\xi-x_0,h)\big|\bigg)\cdot h-
		\bigg(\frac{\vd^2f}{\vd x^2}(x_0,y_0)(\xi-x_0)+\sigma\big|(\xi-x_0,0)\big|\bigg)\cdot h=\\
		=\frac{\vd^2f}{\vd y\cdot\vd x}(x_0,y_0)\cdot h^2+\sigma\big|(\xi-x_0,h)\big|\cdot h+\sigma\big|(\xi-x_0,0)\big|
	\end{multline*}
	Аналогично, только вместо $(\xi)$ - $(\v x_0+h)$, а вместо $(y_0+h)$ - $(\eta)$.\\
	Приравниваем получившиеся выражения:
	\[\frac{\vd^2f}{\vd y\cdot\vd x}(x_0,y_0)+\sigma(1)=\frac{\vd^2f}{\vd x\cdot\vd y}(x_0,y_0)+\sigma(1),\ h\ra0\]
	Отметим, что последний переход имеет место, так как:
	\[\bigg|\frac{\sigma\big|(\xi-x_0,h)\big|\cdot h}{h^2}\bigg|=\bigg|\frac{\sigma(\sqrt{(\xi-x_0)^2+h^2})}{\sqrt{(\xi-x_0)^2+h^2}}\cdot\frac{\sqrt{(\xi-x_0)^2+h^2}}{h}\bigg|=\sigma(1),\ h\ra0\]
\end{proof}
\begin{examples}
	\begin{equation*}
		\begin{split}
			f(x,y)&=
			\begin{cases}
				x^{\frac32}y^{\frac32}\cdot\sin\frac1x\sin\frac1y,\ xy\neq0\\
				0,\ xy=0
			\end{cases}
			\\
			f(x,y)&=
			\begin{cases}
				x^3y^3\cdot\sin\frac{1}{xy},\ xy\neq0\\
				0,\ xy=0
			\end{cases}
			\\
			f(x,y)&=
			\begin{cases}
				x^2\cdot\sin\frac1x,\ x\neq0\\
				0,\ x=0
			\end{cases}
		\end{split}
	\end{equation*}
\end{examples}
\begin{definition}
	Функция называется $m$ раз дифференцируемой в точке $\v x_0$, если все её частные производные порядка $m-1$ определены в окрестности точки $\v x_0$ и дифференцируемы в ней.
\end{definition}
\begin{definition}
	\textit{Дифференциалом второго порядка} дважды дифференцируемой в точке $\v x_0\in\R^n$ функции $f$ называется дифференциал от дифференциала $df$. Отметим, при взятии второго дифференциала дифференциалы независимых переменных те же, что и для дифференциала $df$.
	\[d^2f(\v x_0)=\sum_{i=1}^n\frac{\vd f}{\vd x_i}(\v x_0)dx_i,\ \text{если }x_i\text{ - независимая переменная, то }dx_i=\tr x_i\]
	\begin{multline*}
		d^2f(\v x_0)=d\bigg(\sum_{i=1}^n\frac{\vd f}{\vd x_i}(\v x)dx_i\bigg)(\v x_0)=\sum_{i=1}^nd\bigg(\frac{\vd f}{\vd x_i}\bigg)(\v x_0)\cdot dx_i=\sum_{i=1}^n\sum_{k=1}^n\frac{\vd^2f}{\vd x_k\cdot\vd x_i}(\v x_0)dx_k\cdot dx_i=\\
		=\bigg(\frac{\vd}{\vd x_1}dx_1+\dots+\frac{\vd}{\vd x_n}dx_n\bigg)^2f(\v x_0)=\frac{\vd^2f}{\vd x_1^2}(\v x_0)dx_1^2+2\cdot\frac{\vd^2f}{\vd x_1\cdot\vd x_2}(\v x_0)dx_1\cdot dx_2+\dots
	\end{multline*}
\end{definition}